\subsubsection{Betrag von Summen und Differenzen}
\label{subsubsec:vector_algebra_magnitude_sums_differences}

Nun verwenden wir das Skalarprodukt, um den Betrag einer Summe oder Differenz von
Vektoren zu berechnen. Dies ist beispielsweise relevant, wenn mehrere nicht parallele
Kräfte auf einen Körper wirken und die resultierende Kraft bestimmt werden soll.

\begin{center}
  \CachedTikzFigure{content/resources/chapters/vector_algebra/tikz/vector_sum_difference_magnitudes.tex}
\end{center}

Dazu schreiben wir das Betragsquadrat eines Vektors als sein Skalarprodukt mit
sich selbst und multiplizieren aus:
\begin{ClozeArea}[34]{7}
  \[
    \begin{aligned}
      |\vec a\pm\vec b|^2
      &=(\vec a\pm\vec b)\cdot(\vec a\pm\vec b)\\[0.5em]
      &=\vec a\cdot\vec a+\vec b\cdot\vec b
        \pm2\,\vec a\cdot\vec b\\[0.5em]
      &=|\vec a|^2+|\vec b|^2
        \pm2|\vec a|\,|\vec b|\cos(\varphi).
    \end{aligned}
  \]
\end{ClozeArea}

\begin{KeyPointBox}{Beträge von Vektorkombinationen}{vector_algebra_sum_difference_magnitudes}
  \[
    |\vec a\pm\vec b|
    =\sqrt{|\vec a|^2+|\vec b|^2
      \pm2|\vec a|\,|\vec b|\cos(\varphi)}.
  \]
\end{KeyPointBox}

\begin{KeyPointBox}{Kosinussatz als Differenzfall}{vector_algebra_cosine_rule_derivation}
  Der Kosinussatz ergibt sich unmittelbar aus dem oben betrachteten
  Differenzfall mit $a=|\vec a|$, $b=|\vec b|$ und
  $c=|\vec a-\vec b|$:
  \[
    c^2=a^2+b^2-2ab\cos(\varphi).
  \]
\end{KeyPointBox}

\begin{ExampleBox}{Resultierende Kraft}{vector_algebra_resultant_force_example}
  \noindent
  \begin{minipage}[c]{0.55\linewidth}
    Auf eine Masse wirken zwei Kräfte mit
    \[
      |\textcolor{THMBlue}{\vec F_1}|=\qty{1}{\newton},\qquad
      |\textcolor{olive!80!black}{\vec F_2}|=\qty{2}{\newton}
    \]
    unter einem Winkel von $60^\circ$. Für die resultierende Kraft
    $\vec F_{\mathrm R}=\textcolor{THMBlue}{\vec F_1}
      +\textcolor{olive!80!black}{\vec F_2}$ gilt
    \[
      \begin{aligned}
        |\vec F_{\mathrm R}|
        &=\sqrt{|\textcolor{THMBlue}{\vec F_1}|^2
          +|\textcolor{olive!80!black}{\vec F_2}|^2
          +2|\textcolor{THMBlue}{\vec F_1}|
          \,|\textcolor{olive!80!black}{\vec F_2}|\cos(60^\circ)}\\
        &=\Cloze{4.8cm}{\sqrt7\,\unit{\newton}
          \approx\qty{2.65}{\newton}}.
      \end{aligned}
    \]
  \end{minipage}\hfill
  \begin{minipage}[c]{0.41\linewidth}
    \centering
    \CachedTikzFigure{content/resources/chapters/vector_algebra/tikz/resultant_force_example.tex}
  \end{minipage}
\end{ExampleBox}
